\section{Il Sottobosco: Binario e Numeri (Rif: Slide Lezione 08)}

\section{Stringhe, JSON e Metodi Funzionali (Lez. 08-09)}

\subsection{Metodi per le Stringhe}
Le stringhe in JS sono primitive, ma hanno molti metodi utili per lavorarci.
Ricorda: le stringhe \textbf{non} cambiano mai (sono immutabili). Ogni metodo ne restituisce una nuova.

\begin{concept}{Immutabilità delle Stringhe}
Le stringhe in JS sono primitive e non cambiano mai. I metodi ritornano sempre una **nuova** stringa.
\end{concept}
\begin{jsbox}
let testo = "Ciao Mondo";

// Estrazione
let parte = testo.slice(0, 4); // "Ciao"
let sub = testo.substring(0, 4); // "Ciao"

// Ricerca
let pos = testo.indexOf("Mondo"); // 5

// Trasformazione
let maiusc = testo.toUpperCase(); // "CIAO MONDO"
let pulito = "  ciao  ".trim(); // "ciao" (toglie spazi ai lati)

// Padding (utile per formattare)
let num = "5";
console.log(num.padStart(4, "0")); // "0005"
\end{jsbox}

\subsection{JSON (JavaScript Object Notation)}
\begin{concept}{Serializzazione}
Il JSON è un formato testuale standard. Serve per convertire oggetti in stringhe (per inviarli) e viceversa.
\end{concept}
\begin{jsbox}
let studente = { nome: "Luca", eta: 20 };

// Da Oggetto a Stringa (per inviarlo o salvarlo)
let stringaJSON = JSON.stringify(studente);
console.log(stringaJSON); // '{"nome":"Luca","eta":20}'

// Da Stringa a Oggetto (quando ricevi dati)
let datiRicevuti = '{"nome":"Marco","eta":22}';
let nuovoStudente = JSON.parse(datiRicevuti);
console.log(nuovoStudente.nome); // "Marco"
\end{jsbox}

\subsection{Metodi Avanzati degli Array (Map, Filter, Reduce)}
Invece di usare sempre il ciclo \texttt{for}, JavaScript offre metodi potenti per manipolare gli array.

\subsubsection{1. Map (Trasformare)}
\begin{concept}{Mappatura 1-a-1}
Crea un \textbf{NUOVO} array applicando una funzione a ogni elemento. La lunghezza non cambia.
\end{concept}
\begin{jsbox}
const numeri = [1, 2, 3, 4];
const doppi = numeri.map(x => x * 2);
console.log(doppi); // [2, 4, 6, 8]
\end{jsbox}

\subsubsection{2. Filter (Filtrare)}
\begin{concept}{Selezione Logica}
Crea un \textbf{NUOVO} array con solo gli elementi che superano un test (restituiscono true).
\end{concept}
\begin{jsbox}
const voti = [18, 25, 15, 30, 12];
const promossi = voti.filter(voto => voto >= 18);
console.log(promossi); // [18, 25, 30]
\end{jsbox}

\subsubsection{3. Reduce (Ridurre/Accumulare)}
\begin{concept}{Aggregazione N-a-1}
Riduce tutto l'array a un singolo valore (es. somma, media) usando un accumulatore.
\end{concept}
\begin{jsbox}
const prezzi = [10, 20, 30];
// tot inizia a 0. Per ogni 'p' (prezzo), facciamo tot + p
const totale = prezzi.reduce((tot, p) => tot + p, 0);
console.log(totale); // 60
\end{jsbox}

\noindent\makebox[\linewidth]{\rule{\paperwidth}{0.4pt}}

\subsection{Bitwise Operators: Manipolare i Bit}

\subsubsection*{1. IL CONCETTO}
Possiamo manipolare i singoli bit di un numero. È una roba da bassa lega (driver, protocolli di rete, flag), ma utile per capire come ragiona la macchina.

\subsubsection*{2. GLI OPERATORI}
\begin{itemize}
    \item \texttt{\&} (AND): Entrambi 1.
    \item \texttt{|} (OR): Almeno uno 1.
    \item \texttt{\^} (XOR): Diversi (uno 1 e l'altro 0).
    \item \texttt{\textasciitilde} (NOT): Inverte tutto.
\end{itemize}

\subsubsection*{3. CODICE E ANALISI}
\begin{concept}{Maschere di Bit (Flags)}
Usiamo i bit come interruttori on/off compatti.
\end{concept}
\begin{jsbox}
// Maschere di bit (Flags)
const LETTURA = 1;  // 001
const SCRITTURA = 2; // 010
const ESECUZIONE = 4; // 100

let permessi = LETTURA | SCRITTURA; // 011 (3)

// Controllo: Ho il permesso di lettura?
if (permessi & LETTURA) {
    console.log("Puoi leggere");
}
\end{jsbox}

\subsubsection*{5. TRAPPOLE E ERRORI DI DISTRAZIONE}
\begin{itemize}
    \item Scambiare \texttt{\&} (Bitwise) con \texttt{\&\&} (Logico).
    \begin{itemize}
        \item \texttt{1 \& 2} $\rightarrow$ \texttt{0} (01 \& 10 non hanno bit comuni).
        \item \texttt{1 \&\& 2} $\rightarrow$ \texttt{2} (True \&\& True $\rightarrow$ restituisce l'ultimo valore).
    \end{itemize}
\end{itemize}
